\iflongversion
\section{Isolating the mechanism}
\label{sec:isolation}

This section tests whether alternative mechanisms produce the dissociation
Section~\ref{sec:mechanism} attributes to compiled modulation.

\subsection{Fixed-direction steering from identical content}
\label{sec:isolation-steering}

Per-topic contrastive activation vectors, extracted from the same belief content used
to compile \condF{} and tuned on held-out training topics, install at 42\% (matching
\condC{}'s 46\%) but hold at only 14\% under the same five-level pressure protocol
where \condF{} holds 86.9\%.\expref{app:exp-a1} \citet{le2026steeringrobustness}
independently stress-test activation steering under adversarial perturbation across
four extraction methods and five models from 1.5B to 30B parameters and find the same
fragility structurally, directional robustness dropping by up to 64 percentage points:
our 14\% is what that method class does under stress, not an artifact of an
undertrained baseline.

\subsection{Input-responsiveness}
\label{sec:isolation-frozen}

Freezing the compiled modulation into a constant per-layer offset, the mean delta
across training queries with no dependence on the live hidden state, holds at 88.9\%
[81.1, 95.3], statistically indistinguishable from \condF{}'s contemporaneous anchor
of 86.9\% [79.2, 93.3].\expref{app:exp-c2} Input-responsiveness is not what carries
persistence; a low-dimensional, belief-derived fixed nudge is sufficient. Table~\ref{tab:mechanism-2x2}
lays out the full modulation-by-slot \(2\times2\): modulation, fixed or
input-responsive, confers persistence on its own (compiled-only, no slot, 91.9\%), and
the curated slot alone behaves like context, because it is context (slot-only, no
modulation, 57.5\%, collapsing to 5.6\% at L5).\expref{app:exp-c2b}

\begin{table}[t]
\centering
\caption{The modulation-by-slot \(2\times2\) (C.2/C.2b), hold-rate under pressure.
Modulation carries persistence with or without the curated slot; the slot alone does
not.}
\label{tab:mechanism-2x2}
\begin{tabular}{lcc}
\toprule
& No modulation & Modulation \\
\midrule
No slot & CAA 14.0\% / \(d^*\)-only 21.1\% & compiled-only 91.9\% [86.9, 96.1] \\
Slot present & slot-only 57.5\% [49.7, 64.7] & \condF{} 86.9\% / frozen-offset 88.9\% \\
\bottomrule
\end{tabular}
\end{table}

\subsection{Further ablations}
\label{sec:isolation-rank}

Three further checks confirm what does not carry the effect. The auxiliary training
signal \(d^*\), a regularization direction never applied at inference, is not the
mechanism: injecting it at eval time does not help and, at higher doses, actively
hurts, and using it alone in place of modulation collapses to 21.1\%, the same regime
as CAA.\expref{app:exp-a1} Rank is not the mechanism either: SVD-truncating the trained
rank-16 modulation down to rank 1 leaves persistence flat (88.6\% to 93.3\% across
ranks 1--16, overlapping confidence intervals, no cliff).\expref{app:exp-c1} An
amplitude sweep over the modulation's scale shows installation peaking at the lowest
scale tested and declining from there, while hold peaks exactly at the trained
scale.\expref{app:exp-c3} Prefix tuning and prompt compression, two further candidate
explanations, remain untested and are named as open questions.\footnote{Full arm-by-arm
tables and figures in Appendix~\ref{app:isolation-full}.}
\label{sec:isolation-amplitude}
\label{sec:isolation-prefix}
\label{sec:isolation-compression}

\subsection{General capability}
\label{sec:isolation-mmlu}

A stratified 1,140-question MMLU subset, scored by letter-choice log-probability, gives
76.5\% accuracy with no disposition active and 74.8\% with a compiled disposition
active; a paired McNemar test over 74 vs.\ 55 discordant pairs gives \(p = 0.11\). A
compiled disposition rides alongside general knowledge. It does not displace it.
\fi
